\documentclass[11pt,a4paper]{article}

\usepackage[T1]{fontenc}
\usepackage[utf8]{inputenc}
\usepackage{amsmath,amssymb}
\usepackage{array}
\usepackage{listings}
\usepackage{url}
\newcommand{\href}[2]{#2}
\Urlmuskip=0mu plus 1mu
\emergencystretch=2em
\widowpenalty=10000
\clubpenalty=10000
\displaywidowpenalty=10000

\newtheorem{definition}{Definition}
\newtheorem{proposition}{Proposition}
\newtheorem{theorem}{Theorem}
\newtheorem{remark}{Remark}
\newenvironment{proof}{\par\noindent\textit{Proof.}\ }{\hfill$\square$\par}

\newcommand{\Types}{\mathcal{T}}
\newcommand{\Symbols}{\Sigma}
\newcommand{\Effects}{\mathcal{E}}
\newcommand{\meet}{\mathbin{\sqcap}}
\newcommand{\seq}{\mathbin{\cdot}}
\newcommand{\zero}{\mathbf{0}}
\newcommand{\id}{\mathbf{1}}
\newcommand{\compact}{\mathsf{compact}}
\newcommand{\eval}[1]{[\![#1]\!]}

\lstset{
  basicstyle=\ttfamily\small,
  columns=fullflexible,
  keepspaces=true,
  showstringspaces=false,
  frame=single,
  breaklines=true
}

\title{Configurable Source-Level Type Checking for Forth\\with Symbolic Stack Effects}
\author{Jaanus Pöial\\Tallinn University of Technology}
\date{}

\begin{document}
\maketitle

\begin{abstract}
Forth words communicate through an implicit data stack, and Forth source can
also change the parser and dictionary while it is being read.  This paper
describes a configurable, source-level checker for that setting.  A phrase is
represented by an input and an output stack of typed symbolic cells.  Repeated
indices record that two positions contain the same abstract value.  Linear
phrases are checked by unifying the touching stack boundaries.  At a branch,
the checker must replace two path effects by one interface; it does so with a
partial syntactic common-contract operation that also handles unequal visible
depths when the difference is an unchanged stack prefix.  Repetition is
handled by comparing one execution with two, not by computing a general loop
invariant.

The type hierarchy, primitive effects, parser words, defining words, literals,
and structured control forms are supplied in text files.  The paper gives
step-by-step definitions of the three effect operators, a complete runnable
example, an artifact manifest, and an exact characterization of the
bottom-up computation performed after a body has been parsed.  It also
separates guarantees from implementation choices: normalization can discard
identity information, branch combination has no accompanying semantic
soundness theorem, and a concrete counterexample shows that the one-versus-two
loop test does not establish correctness for arbitrary iteration counts.  The
native implementation is \texttt{gforth-evaluator.fs}; comparisons with
StrongForth, an optional pluggable Forth type system, and WebAssembly explain
which limitations arise from this checker rather than from stack languages in
general.
\end{abstract}

\noindent\textbf{Keywords:} concatenative language; Forth; static analysis;
stack effect; abstract interpretation; symbolic execution.

\section{Introduction}

\par\noindent\begin{minipage}{\linewidth}
In a concatenative language, juxtaposition is composition: if $p$ and $q$ are
programs, then $p\,q$ executes $p$ and passes its resulting stack directly to
$q$.  A conventional Forth stack annotation has the form
\begin{lstlisting}
( n n -- n )
\end{lstlisting}
\end{minipage}\par
Such an annotation resembles a type signature, but an ordinary comment has no
static force and does not generally express correlations between stack cells.
For example,
\par\begin{samepage}
\noindent\begin{minipage}{\linewidth}
\begin{lstlisting}
DUP  ( x[1] -- x[1] x[1] )
OVER ( x[2] x[1] -- x[2] x[1] x[2] )
\end{lstlisting}
\end{minipage}
\end{samepage}\par
records facts that the unindexed signatures cannot express: two output
positions may denote the same symbolic input.

This project turns such signatures into an executable abstract semantics.  It
loads three inputs: a finite type hierarchy, a vocabulary of word and parser
specifications, and source text.  It infers effects for sequences and newly
defined words without executing the concrete program.  The implementation is
a single native Gforth source file, \texttt{gforth-evaluator.fs}.

For readers unfamiliar with Forth, a \emph{word} is a named operation and a
colon definition \texttt{: NAME ... ;} introduces a new word.  Operands are
normally taken from the data stack rather than from named arguments; throughout
this paper the top cell is the rightmost one.  Thus \texttt{1 2 +} pushes two
integer literals and then replaces them by their sum.  Words such as
\texttt{DUP}, \texttt{SWAP}, and \texttt{OVER} respectively copy the top cell,
exchange the top two cells, and copy the second cell over the top one.  Control
words such as \texttt{IF}, \texttt{ELSE}, and \texttt{THEN} execute while a
definition is compiled and delimit source regions that run later.  These data-
and compile-time roles are kept separate below.

The analysis addresses four problems in addition to stack-height consistency.
\begin{enumerate}
  \item compatibility of an effect's output boundary with the next effect's
        input boundary;
  \item propagation of subtype refinement and symbolic identity through stack
        shuffles;
  \item construction of a common contract for alternative branches;
  \item finite approximation of unbounded iteration.
\end{enumerate}

The answers implemented here are, respectively, comparable-type unification,
sequence-wide substitution, a partial effect meet, and a local repeated-effect
operator.  These are choices made by this checker, not limitations inherent to
concatenative languages.  For example, WebAssembly validates loops against
declared control-label types, and other typed Forth systems make different
choices about annotations and overloading (Section~\ref{sec:related}).  The
present choices continue the algebraic development of stack effects in
\cite{poial1991,poial1992,poial1994,poial2002,poial2006}.  An earlier Java
implementation of the framework was published in 2008 \cite{poial2008}; the
implementation studied in this paper is the native Gforth version.

The analysis distinguishes five technical properties.
\begin{itemize}
  \item The declared type relation need not be a complete lattice.  The core
        algorithm requires only a finite normalized relation and selects a
        meet when the two types are comparable.
  \item Raw effects and printed effects are different domains.  Normalization
        is deterministic, but it may erase an implicitly introduced
        two-occurrence correlation.
  \item Branch merging is an operation on contracts, not the set-theoretic
        union of branch transition relations.
  \item The implemented loop operation is $\pi\meet(\pi\seq\pi)$.  It is not a
        proof that the result describes every iteration count.
  \item Once parsing and leaf lookup are fixed, checking is the unique
        bottom-up evaluation of an acyclic system of effect equations.
\end{itemize}

\section{Type and Symbol Domains}

\subsection{Normalized finite subtype relation}

Let $\Types$ be the finite set of canonical type names.  The input permits
aliases; mathematically, aliases are quotiented before considering the strict
subtype declarations.  Write $\tau\leq\upsilon$ when $\tau$ is at least as
specific as $\upsilon$.  The bundled Forth-like profiles contain chains such
as
\[
  \mathit{a\mbox{-}addr} < \mathit{c\mbox{-}addr}
  < \mathit{addr} < u < x,
  \qquad
  \mathit{char} < {+}n < n < x.
\]
The implementation adds identity, inverse lookup information, and transitive
relations after loading the declarations.

\begin{definition}[Comparable type meet]
For $\tau,\upsilon\in\Types$, define the partial operation
\[
\tau\meet\upsilon =
\begin{cases}
 \tau, & \tau\leq\upsilon,\\
 \upsilon, & \upsilon\leq\tau,\\
 \text{undefined}, & \text{otherwise}.
\end{cases}
\]
\end{definition}

Thus the implementation does not search for an arbitrary third subtype when
$\tau$ and $\upsilon$ are incomparable.  This is stricter than meet in a
general lattice and identifies failures at the corresponding boundary.

\subsection{Typed symbolic cells}

A stack symbol is a triple
\[
  a=(\tau,i,e)\in\Types\times\mathbb{N}\times\{0,1\}.
\]
It is printed as $\tau$ when $i=0$ and as $\tau[i]$ otherwise.  The bit $e$
records whether an index was explicit in the specification.  Symbols with the
same positive index are correlated.  An unindexed occurrence is freshened
before composition, so two occurrences of plain \texttt{x} do not accidentally
become equal.

This system uses indices as first-order equality variables, not as names of
runtime storage locations.  Equal indices assert one correlation; different
indices do not assert that the concrete values are unequal.  A substitution
$[a\mapsto b]$ replaces every occurrence of $a$ in the current effect sequence.
Global replacement is what allows a type fact discovered at one composition
boundary to propagate through \texttt{DUP}, \texttt{SWAP}, or \texttt{OVER}.

The brackets are an internal, explicit version of the Forth standard's numeric
suffix convention.  Standard prose commonly writes \texttt{n1 n2 -- n3} to
distinguish several items of type \texttt{n}, while a word such as
\texttt{DUP ( x -- x x )} needs no suffix because every shown \texttt{x}
refers to the copied item.  In formulas below, \texttt{x[1]} makes the equality
variable unambiguous.  When effects are presented to a Forth user, indices
should be renumbered consecutively and omitted only where the standard
convention remains unambiguous.

\section{Stack Effects and Frames}

Let $\Symbols$ be the set of typed symbols and let $\Symbols^*$ be finite
symbol sequences.  The top of stack is written on the right.

\begin{definition}[Raw effect]
A raw stack effect is a pair
\[
  \pi=(L\;--\;R),\qquad L,R\in\Symbols^*.
\]
It denotes a transformation of a visible stack suffix: for an arbitrary
unchanged frame $F$, a stack of shape $F L$ is transformed to $F R$, subject
to the type and equality constraints represented by the symbols.
\end{definition}

The empty effect $\id=(\;--\;)$ is identity.  Algorithmic failure is represented
by a separate clash value $\zero$, not by an ordinary effect.

The implicit frame is essential.  For example, $(n--n)$ does not assert that
the whole runtime stack has exactly one cell; it constrains only the suffix
touched by the word.  Effects with different visible depths may consequently
still be combined if their depth difference is a preserved prefix.  Type-system
literature often calls such an arbitrary sequence a \emph{stack row} and writes
a row variable, say $\rho$, for it.  The effect above could therefore be read
$\rho\,n\to\rho\,n$.  A row variable ranges over whole stack prefixes, whereas
the indexed variables of Section~2 range over single cells.

For raw effects, alpha-renaming positive indices has no meaning.  Let
$\pi\equiv_\alpha\pi'$ mean equality up to a bijective renaming of indices that
preserves types and explicitness.  Implementations must freshen independently
instantiated word specifications before attempting unification.

\section{Sequential Composition}
\label{sec:composition}

Let
\[
  \pi_1=(L_1--R_1),\qquad \pi_2=(L_2--R_2).
\]
Composition inspects the touching boundary from its top inward.  If the current
top symbols are $a=\tau[i]$ and $b=\upsilon[j]$, calculate
$\rho=\tau\meet\upsilon$.  If undefined, composition clashes.  Otherwise make
a fresh symbol $z=\rho[k]$, substitute $a$ and $b$ by $z$ everywhere in both
working effects, remove the matched output/input pair, and continue.

A direct functional presentation follows.  All indices in the two operands
are assumed to be disjoint.  Let $\mathcal{C}$ denote raw composition.  The
first two cases are
\[
\begin{aligned}
\mathcal{C}(L_1,\epsilon,L_2,R_2)&=(L_2L_1--R_2),\\
\mathcal{C}(L_1,R_1,\epsilon,R_2)&=(L_1--R_1R_2).
\end{aligned}
\]
For the remaining case write $R_1=\widehat R_1a$ and
$L_2=\widehat L_2b$.  If the types of $a$ and $b$ are incomparable,
$\mathcal{C}$ is undefined.  Otherwise let $z$ be their fresh, more specific
representative and let $\sigma=[a\mapsto z,b\mapsto z]$.  Then
\[
 \mathcal{C}(L_1,R_1,L_2,R_2)=
 \mathcal{C}(\sigma L_1,\sigma\widehat R_1,
             \sigma\widehat L_2,\sigma R_2).
\]
This notation makes every changed sequence explicit.  In the implementation,
substitution mutates all four working sequences first; after that mutation the
same recursive call is written
$\mathcal{C}(L_1,R'_1,L'_2,R_2)$.  In particular, no unexplained
$L'_1$ or $R'_2$ is required.  Concatenation order in the first base case
reflects the implicit frame: inputs still required by $\pi_2$ lie beneath
those already required by $\pi_1$.

Define the totalized product
\[
\begin{aligned}
\pi_1\odot\pi_2&=
\begin{cases}
 \compact(\pi_1\seq\pi_2),&\text{if raw composition succeeds},\\
 \zero,&\text{otherwise},
\end{cases}\\
\zero\odot\pi&=\pi\odot\zero=\zero.
\end{aligned}
\]
This is the operation used while reducing a source sequence.

\paragraph{Worked derivation.}
With
\par\begin{samepage}
\noindent\begin{minipage}{\linewidth}
\begin{lstlisting}
1    literal integer ( -- n )
2    literal integer ( -- n )
+                    ( n n -- n )
DUP                  ( x[1] -- x[1] x[1] )
\end{lstlisting}
\end{minipage}
\end{samepage}\par
each literal introduces a fresh numeric cell.  The fragment
\texttt{1 2 + DUP} then has the derivation
\[
(\!--n)\seq(\!--n)=(\!--n\;n),
\]
\[
(\!--n\;n)\seq(n\;n--n)=(\!--n),
\]
\[
(\!--n)\seq(x[1]--x[1]\;x[1])
=(\!--n[1]\;n[1]).
\]
The last step refines \texttt{x} to \texttt{n} and propagates that refinement
to both outputs of \texttt{DUP}.  Running this example through the native
evaluator with \texttt{ex1types.txt} and \texttt{ex1specs.txt} confirms the
calculation:
\par\begin{samepage}
\noindent\begin{minipage}{\linewidth}
\begin{lstlisting}
1    ( -- n )
2    ( -- n )
+    ( n n -- n[1] )
DUP  ( n[1] -- n[1] n[1] )
< n[1] n[1]
\end{lstlisting}
\end{minipage}
\end{samepage}\par

\begin{proposition}[Determinacy of raw composition]
For fixed freshening and fresh-index policies, raw boundary composition either
returns one effect or clashes.  Different choices of fresh positive indices
produce alpha-equivalent effects.
\end{proposition}

\begin{proof}
At each recursive step there is one pair of top boundary symbols.  The
comparable meet, when it exists, is unique.  Substitution and removal therefore
determine the next state.  The recursion decreases
$\lvert R_1\rvert+\lvert L_2\rvert$ by two and terminates.  Fresh-index choices
can change only variable names, hence the alpha-equivalence result.
\end{proof}

This proposition applies to raw effects.  Algebraic associativity of an ideal
stack-effect product motivates the polycyclic-monoid account in
\cite{poial1991,poial1994}; however, inserting a lossy compaction after every
binary product can affect later identity information.  Consequently, algebraic
laws of the raw domain do not automatically hold for the stored
representation.

\section{Internal Identities and Forth-Style Output}

There are two separate normalization tasks.  Internally, alpha-equivalent
identity variables should be renamed deterministically but retained.  For
presentation, the Forth standard's style numbers same-typed items only as
needed and renumbers them consecutively.  For example, an internal effect
\[
 (x[7]\;x[8]--x[7]\;x[9])
\]
can be displayed as \texttt{( x1 x2 -- x1 x3 )}: the first result is the first
input, while the other result has no stated correlation.  The special case
\texttt{DUP ( x -- x x )} needs no suffix because all displayed \texttt{x}
items denote the same value.

The current Gforth artifact does less.  After a successful operation it counts
symbol occurrences over the result and its working sequence, traverses them in
a deterministic order, and assigns small indices.  An index remains visible
if it was explicit in a source specification or if the identity occurs more
than twice.  An automatically introduced identity occurring once or twice is
set to zero and printed without an index.  Thus an implicit
$(n[7]--n[7])$ becomes $(n--n)$.  More seriously, the three identities in the
example above can be printed as \texttt{( x x -- x x )}, hiding which input is
preserved.

The only practical advantage of that policy is short, backward-compatible
output.  It gives the analysis no additional power.  Standard-style numbering
is preferable for inferred effects: retain every equivalence class internally,
renumber classes consecutively for display, and omit suffixes only when doing
so is unambiguous under the Forth convention.

\begin{remark}[Current compaction is a widening]
Removing an index preserves stack shape and per-position type constraints but
forgets an equality constraint.  Consequently the implemented $\compact$ is
deterministic and concise, but neither injective nor semantics-preserving for
symbolic identity.  It is a conservative loss of correlation information.
\end{remark}

This is not merely a printer issue in the current artifact: compacted effects
are stored and reused.  A faithful formalization must therefore define the
implemented operators as ``raw operation followed by compaction.''  A revised
implementation should keep raw identities in the dictionary and apply
Forth-style formatting only at the output boundary.

\section{Branches: Two Paths, One Interface}
\label{sec:branches}

The central branch problem arises after, not inside, an \texttt{IF}.  The two
arms may require different visible suffixes, leave different cells, or preserve
different identities, but the following word and the inferred definition need
one interface.  Equal stack depth is necessary but not sufficient: for
example, one arm may leave an address where the other leaves a number.  A
checker must either retain a disjunction of path effects, forget enough facts
to form a safe common summary, restrict the accepted programs, or combine
these strategies.  This checker chooses one partial, syntactic contract.

The operation is named \texttt{glb} in the code and is written $\meet$ here.
The symbol $\sqcap$ does mean ``meet.''  It records the intended
order-theoretic reading---greatest lower bound in a contract order---but the
actual operator is the following algorithm, not an appeal to an unspecified
lattice.

Let $\pi=(L--R)$ and $\theta=(M--N)$.  Suppose first that $|L|>|M|$ and put
$k=|L|-|M|$.  A common effect is possible only if
\[
                         |R|-|N|=k.                 \tag{1}
\]
Write the longer sides as
\[
                L=P\,L_s,\qquad R=Q\,R_s,
                \qquad |P|=|Q|=k.
\]
The deep prefixes $P$ and $Q$ are the part of the implicit frame made visible
by the longer effect.  The algorithm first unifies $P$ with $Q$ pairwise, so
that this prefix is preserved.  It then unifies $L_s$ with $M$ and $R_s$ with
$N$ pointwise.  The case $|M|>|L|$ is symmetric, and equal input lengths are
the special case $k=0$.  A length failure in (1), or any incomparable pair of
types, makes the operation undefined.

Here is the unequal-length case explicitly:
\[
 \pi=(x[1]\;n--x[1]\;\mathit{char}),\qquad
 \theta=(n--\mathit{char}).
\]
For $\pi$, one cell of the otherwise implicit frame is displayed.  We have
$k=1$ on both sides; $x[1]$ is unchanged, and the residual
$n--\mathit{char}$ agrees with $\theta$.  Their common contract is therefore
$\pi$.  In contrast,
\[
 (x--x\;x)\qquad\hbox{and}\qquad(\epsilon--\epsilon)
\]
have input-length difference one and output-length difference two, so they do
not merge.  This is why \texttt{IF DUP THEN}, with no corresponding
\texttt{ELSE}, is rejected: only one run-time path leaves the extra cell.

With equal lengths, type refinement is pointwise.  For example,
\[
  (x\;\mathit{a\mbox{-}addr}--\epsilon),\qquad
  (x\;\mathit{c\mbox{-}addr}--\epsilon)
\]
produce the stronger requirement
$(x\;\mathit{a\mbox{-}addr}--\epsilon)$.  Composing the declared
\texttt{IF ( flag -- )} effect then accounts for the selector.

The unequal-depth rule above also exposes an implementation detail.  The
mathematical alignment compares the residual slices at offset $k$.  The
current native helper \texttt{ev-spec-unify} starts its final vector traversal
at offset zero.  Consequently the artifact can reject an unequal-depth merge
when the displayed frame prefix and residual suffix have different types,
even though the rule above accepts it.  Cases in which broad \texttt{x} cells
make both alignments compatible can conceal this discrepancy.  It is recorded
as a current implementation limitation rather than built into the rule.

A second limitation appears even without type differences.  Consider
\begin{lstlisting}
: MAYSWAP  IF SWAP THEN ;
\end{lstlisting}
The true arm swaps two cells and the false arm preserves them.  A conventional
single Forth summary that makes no unsupported value claim is
\begin{lstlisting}
( x1 x2 flag -- x3 x4 )
\end{lstlisting}
while a precise analysis would retain the disjunction ``\texttt{x2 x1} or
\texttt{x1 x2}.''  The current syntactic meet instead makes the two inputs one
identity and reports
\begin{lstlisting}
( x[1] x[1] flag -- x[1] x[1] )
\end{lstlisting}
That restricted contract is not the general effect of \texttt{MAYSWAP}.  It
illustrates why a syntactic meet that strengthens preconditions must not be
confused with a sound over-approximation of both run-time paths.

Thus $\meet$ is useful operational notation, but the name \texttt{glb} alone
does not prove that the code computes a semantic greatest lower bound.  Such a
proof needs an operational meaning for primitive effects, an explicit
variance-aware contract order, and a decision about whether path summaries may
be disjunctive.  No such semantic soundness theorem is claimed here.

\section{Iteration: What One Versus Two Can Establish}
\label{sec:loops}

A loop has a back edge: the abstract stack produced by one pass must be usable
again at the loop header.  A conventional abstract interpreter records a
header state $I$, propagates the body through it, joins the result back into
$I$, and iterates until a post-fixpoint is reached (with widening when needed)
\cite{cousot1977}.  WebAssembly takes another route: block and loop label types
are declared, and validation checks every transfer against those types.  The
present checker does neither.  It uses a local operation inherited from the
stack-effect calculus.

\begin{definition}[One-versus-two repeated summary]
For an effect $\pi$, define
\[
  \pi^\star_2 = \pi\meet(\pi\seq\pi)
\]
when both raw composition and the syntactic effect meet succeed; compact the
result before storing it.
\end{definition}

The test is finite and catches any nonzero net stack-depth change, because the
depth changes of $\pi$ and $\pi^2$ are respectively $d$ and $2d$.  Agreement
requires $d=0$.  That useful depth fact must not be generalized into a claim
that the returned symbolic effect is an invariant for every iteration count.
The problem can be seen without subtyping.

Let \texttt{x} be the only data-cell type and declare two synthetic
primitives; \texttt{NEXT?} obtains an unconstrained termination flag:
\begin{lstlisting}
P      ( x[1] x[2] -- x[3] x[1] )
NEXT?  ( -- flag )
UNTIL  ( flag -- )
\end{lstlisting}
The last two effects cancel at the data boundary, so one complete loop pass has
the effect of \texttt{P}.  On concrete symbolic names, that pass maps
$A\,B$ to $C\,A$, two passes map it
to $D\,C$, and three map it to $E\,D$.  Up to consecutive renaming,
\[
\begin{aligned}
 \pi   &=(x[1]\;x[2]--x[3]\;x[1]),\\
 \pi^2 &=(x[1]\;x[2]--x[3]\;x[4]),\\
 \pi^3 &=(x[1]\;x[2]--x[3]\;x[4]).
\end{aligned}
\]
Only the one-pass effect says that the top result is the original lower input.
Nevertheless, the current syntactic meet can unify the unrelated variable in
$\pi^2$ with that input.  Because \texttt{NEXT?} may return true after any pass, this loop has arbitrary
finite execution counts, yet it is accepted:
\begin{lstlisting}
: TEST  BEGIN P NEXT? UNTIL ;
\end{lstlisting}
and the native artifact reports
\begin{lstlisting}
TEST ( x[2] x[1] -- x[3] x[2] )
\end{lstlisting}
which is alpha-equivalent to the one-pass correlation.  Composing that reported
summary with itself produces the uncorrelated two-pass effect, not the reported
summary.  Therefore a checker requiring an inductive or idempotent summary
would reject this candidate (or weaken it to an uncorrelated summary) even
though the one-versus-two routine accepts it.

This example also qualifies the phrase ``strictly weaker.''  In a depth-only
domain the one-versus-two test already proves zero net depth, and a fixpoint
analysis may accept the same loop with a less precise invariant.  The precise
claim is not that every fixpoint checker rejects more programs.  It is that
$\pi\meet\pi^2$, with the implemented syntactic meet, is not by itself a proof
that its returned type-and-identity contract describes all $\pi^j$.

The native word \texttt{ev-spec-pistar} implements exactly this local rule, and
the declarative \texttt{repeat} constructor calls it.  It does not calculate a
control-flow fixpoint or check the result for inductiveness.  The notation
$\pi^\star_2$ is used to avoid confusing the implementation with Kleene star.

\section{Declarative Source Semantics}

The effect algebra is only one layer of a source checker.  Forth parser words
can consume names or delimited text, definitions extend the dictionary, and
control words are immediate.  The specification format attaches metadata such
as
\par\begin{samepage}
\noindent\begin{minipage}{\linewidth}
\begin{lstlisting}
"("      parse until ")"       ( -- )
CONSTANT parse word define     ( x -- )
VARIABLE parse word define     ( -- a-addr )
:        parse word define colon ( -- )
IF       control if            ( flag -- )
\end{lstlisting}
\end{minipage}
\end{samepage}\par
Ordinary consumed comments are not trusted as type declarations.  Thus a
source comment \texttt{( -- n )} cannot override inference.  Recognized
locals syntax may provide a provisional definition shape.

There are two different effects around a control word.  For example,
\texttt{IF} consumes a run-time data-stack flag, but while compiling it also
places an \texttt{orig} item on the Forth control-flow stack; \texttt{THEN}
resolves that item \cite{forth2012}.  The present profile records the first as
\texttt{( flag -- )} and records the second only as parser metadata such as
\texttt{control if} or \texttt{control fi}.  Structural parsing catches
misnested or missing delimiters, but the control-flow stack is not represented
as a second typed stack effect.  Adding effects over standard items such as
\texttt{orig}, \texttt{dest}, and \texttt{do-sys} would permit compile-time
stack checking independently of the run-time data-stack analysis.

Structured controls are supplied through paired \texttt{syntax:} and
\texttt{effect:} descriptions.  A conditional can be specified as follows (the demo spelling \texttt{FI} is a
synonym for the usual \texttt{THEN}):
\par\begin{samepage}
\noindent\begin{minipage}{\linewidth}
\begin{lstlisting}
syntax:
  IF <then branch> [ELSE <else branch>] FI
effect:
  IF
  either <then branch> <else branch>
\end{lstlisting}
\end{minipage}
\end{samepage}\par
The effect-template language is
\[
E ::= \epsilon \mid \langle s\rangle \mid c \mid E_1E_2
    \mid \mathsf{either}(E_1,E_2) \mid \mathsf{repeat}(E),
\]
where $\langle s\rangle$ refers to a captured source segment and $c$ is a
control-word runtime effect.  Its implemented semantics is
\begin{align*}
\eval{\epsilon}&=\id,\\
\eval{\langle s\rangle}&=X_s,\\
\eval{c}&=\Gamma(c),\\
\eval{E_1E_2}&=\eval{E_1}\odot\eval{E_2},\\
\eval{\mathsf{either}(E_1,E_2)}&=
  \overline{\meet}(\eval{E_1},\eval{E_2}),\\
\eval{\mathsf{repeat}(E)}&=
  \overline{\star}_2(\eval{E}),
\end{align*}
where bars denote totalization with $\zero$ and include compaction after a
successful raw operation.

\section{The \texttt{ex1} Input/Output Example}
\label{sec:ex1}

The bundled \texttt{real} profile gives a complete, reproducible example of
the evaluator interface.  It consists of the three files
\texttt{ex1types.txt}, \texttt{ex1specs.txt}, and \texttt{ex1prog.txt}.  The
launcher command is
\par\begin{samepage}
\noindent\begin{minipage}{\linewidth}
\begin{lstlisting}
./run-evaluator-gforth.sh real
\end{lstlisting}
\end{minipage}
\end{samepage}\par
or, equivalently, an explicit invocation of
\texttt{gforth-evaluator.fs} with \texttt{--types}, \texttt{--specs}, and
\texttt{--prog} arguments.

\par\noindent\begin{minipage}{\linewidth}
\subsection{Type input}

The type file declares canonical names, aliases, subtype edges, and scanner
delimiters.  Its essential declarations are:
\begin{lstlisting}
type X cell x
type n number N
type char c
type flag boolean f
type addr a
type c-addr ca
type a-addr aa
type xt

rel n < X
rel flag < n
rel addr < X
rel char < n
rel c-addr < addr
rel a-addr < c-addr
rel xt < X
scanner EOL "\n"
\end{lstlisting}
\end{minipage}\par
For example, \texttt{X}, \texttt{cell}, and \texttt{x} name one canonical
type, while \texttt{a-addr} is more precise than \texttt{c-addr}.  The
loader computes the transitive closure, so it also knows
$\mathit{a\mbox{-}addr}<\mathit{addr}<X$.

\subsection{Specification input}

The specification file supplies both runtime effects and source-processing
metadata.  Representative entries are:
\par\begin{samepage}
\noindent\begin{minipage}{\linewidth}
\begin{lstlisting}
DUP      ( X[1] -- X[1] X[1] )
SWAP     ( X[2] X[1] -- X[1] X[2] )
"("      parse until ")" ( -- )
:        parse word define colon ( -- )
;        control end ( -- )
literal integer ( -- n )
IF       control if ( flag -- )
ELSE     control else ( -- )
FI       control fi ( -- )
BEGIN    control begin ( -- )
WHILE    control while ( flag -- )
REPEAT   control repeat ( -- )
UNTIL    control until ( flag -- )
DP       ( -- a-addr )
@        ( a-addr -- X )
C@       ( c-addr -- char )
0=       ( n -- flag )
\end{lstlisting}
\end{minipage}
\end{samepage}\par
Thus the same input describes ordinary words, literal recognition, parser
words, defining words, and control roles.  The indexed effects of \texttt{DUP}
and \texttt{SWAP} preserve correlations; \texttt{@} and \texttt{C@} exercise
the address subtype chain.  \texttt{DP} is a system-specific demo word denoting
the data-space (dictionary) pointer; it is not assumed implicitly and has the
declaration \texttt{DP ( -- a-addr )} shown above.

The current dictionary accepts one effect declaration per word, so it cannot
directly overload \texttt{0=} with separate \texttt{( n -- flag )} and
\texttt{( u -- flag )} entries.  A profile can instead declare their common
accepted supertype.  For maximum Forth compatibility, the bundled
\texttt{forth2012specs.txt} uses \texttt{0= ( x -- flag )}; a stricter numeric
profile can introduce an \texttt{integer} type above both \texttt{n} and
\texttt{u} and use \texttt{( integer -- flag )}.  True multiple-effect
selection, as supported in other Forth typing work, remains a useful extension.

\subsection{Program input and inferred output}

The program file defines five words and then checks one top-level sequence:
\par\begin{samepage}
\noindent\begin{minipage}{\linewidth}
\begin{lstlisting}
: PEEK2 DUP @ SWAP @ ;
: MAYSWAP IF SWAP FI ;
: NLOOP BEGIN DUP 0= WHILE REPEAT ;
: ULOOP BEGIN DUP 0= UNTIL ;
: KEEPIF IF DUP ELSE DUP FI ;

DP DP C@ 0= KEEPIF
\end{lstlisting}
\end{minipage}
\end{samepage}\par
The source layout is immaterial outside parsed strings and comments.  The
native run writes the following exact effects to \texttt{ex1prog.txt.log}:
\par\begin{samepage}
\noindent\begin{minipage}{\linewidth}
\begin{lstlisting}
PEEK2   ( a-addr[1] -- X[2] X )
MAYSWAP ( X[1] X[1] flag -- X[1] X[1] )
NLOOP   ( n[1] -- n[1] )
ULOOP   ( n[1] -- n[1] )
KEEPIF  ( X[1] flag -- X[1] X[1] )
\end{lstlisting}
\end{minipage}
\end{samepage}\par
The \texttt{MAYSWAP} line is reproducible artifact output, but it is not the
proper general interface of that word.  In standard-style notation the safe
non-disjunctive summary is \texttt{( X1 X2 flag -- X3 X4 )}; a more precise
system would retain the two possible output orders.  The repeated index in the
actual line makes every \texttt{X} one symbolic value and is too restrictive,
as explained in Section~\ref{sec:branches}.

The top-level trace ends with
\par\begin{samepage}
\noindent\begin{minipage}{\linewidth}
\begin{lstlisting}
DP      ( -- a-addr[1] )
DP      ( -- a-addr )
C@      ( a-addr -- char )
0=      ( char -- flag )
KEEPIF  ( a-addr[1] flag -- a-addr[1] a-addr[1] )
< a-addr[1] a-addr[1]
\end{lstlisting}
\end{minipage}
\end{samepage}\par
Type refinement occurs at \texttt{C@}: the producer yields
\texttt{a-addr}, while \texttt{C@} requires \texttt{c-addr}.  Since
$\mathit{a\mbox{-}addr}<\mathit{c\mbox{-}addr}$, composition succeeds and
retains the more precise address information.  The flag from \texttt{0=} is
then consumed by \texttt{KEEPIF}; both branches duplicate the surviving
address, yielding the final pair of correlated \texttt{a-addr} cells.  The
current native evaluator was run on both this complete profile and the
command-line example in Section~\ref{sec:composition}; both produced the
outputs shown above.

\subsection{Acyclic effect equations}

Fix a finite parsed tree $T$ for one source body and environment $\Gamma$.
Assume every undecomposed leaf has already been assigned an effect by word
lookup, literal classification, or an earlier definition.  Associate a
variable $X_A$ with every tree node $A$:
\begin{align*}
 X_A&=\Gamma(A) &&\text{if $A$ is a leaf},\\
 X_A&=X_B\odot X_C &&\text{if $A$ sequences children $B,C$},\\
 X_A&=\widehat{\eval{E}}_X &&\text{if $A$ applies control template $E$}.
\end{align*}
Every right-hand side mentions only proper descendants, so the dependency
graph is acyclic.

\begin{theorem}[Exact characterization of parsed-body checking]
Let $T$ and $\Gamma$ satisfy the assumptions above.  The induced equation
system has exactly one solution in
\[
D=\{\text{compacted effects}\}\cup\{\zero\}.
\]
For every segment $A$, the body evaluator returns its solution value $X_A$;
it reports a clash at $A$ exactly when $X_A=\zero$.  The body is accepted if
and only if the root value is nonzero.
\end{theorem}

\begin{proof}
Proceed by induction on node height.  A leaf is fixed by $\Gamma$.  At a
sequence node, the induction hypothesis uniquely fixes both child values and
$\odot$ is a deterministic total function on $D$.  At a control node,
structural induction on $E$ uniquely determines the result because each
constructor invokes a deterministic totalized operator.  The evaluator visits
the same children and applies the same operator.  Therefore it computes the
unique equation value at every node, including the root.
\end{proof}

This is an exactness theorem for the implemented, post-parse computation; it is
not a conventional progress-and-preservation theorem.  It does not cover the
outer interpreter's evolving dictionary, provisional forward-reference
summaries, recursion, or the semantic correctness of primitive specifications.

\section{Native Gforth Implementation}

The implementation is the single source file
\texttt{gforth-evaluator.fs}.  Its word families mirror the formal layers:
\begin{center}
\begin{tabular}{@{}p{0.39\linewidth}p{0.55\linewidth}@{}}
\hline
Formal role & Representative Gforth words\\
\hline
canonical types and relation & \texttt{ev-ts-*}\\
typed indexed symbols & \texttt{ev-sym-*}\\
stack-effect records & \texttt{ev-spec-*}\\
linear reduction and normalization & \texttt{ev-spec-list-*}\\
specification dictionary & \texttt{ev-ss-*}\\
source scanning and program parsing & \texttt{ev-sc-*}, \texttt{ev-parse-*}\\
diagnostics and log output & \texttt{ev-diag-*}, \texttt{ev-log-*}\\
command-line driver & \texttt{ev-parse-args}, \texttt{ev-run-native}\\
\hline
\end{tabular}
\end{center}

Strings, growable vectors, source spans, type symbols, effects, and control
expressions are represented by explicitly allocated Gforth records.  The
checker normalizes CR, LF, and CRLF input, clones mutable dictionary effects
before evaluation, allocates disjoint wildcard ranges, and builds the final
annotation from the same normalized working value as the final effect.  The
launcher reserves a 4 MiB dictionary.

For effects of visible sizes $m$ and $n$, boundary matching itself performs at
most $\min(m,n)$ recursive matches.  The actual cost is higher because each
match substitutes through the current working vectors.  With this
representation, reducing a program of $N$ words is roughly quadratic in the
largest intermediate symbolic effect in the worst case.  A union--find
representation for identities and persistent stack rows would reduce copying.

Mutation imposes three implementation obligations:
\begin{enumerate}
  \item dictionary effects must be deep-cloned before evaluation;
  \item independently evaluated branches need disjoint wildcard ranges;
  \item normalization and wildcard allocation must be deterministic across
        repeated native runs.
\end{enumerate}
All three are handled explicitly by the current Gforth source.

\section{Artifact Contents and Evaluation}

\subsection{Manifest and reproduction}

The example in Section~\ref{sec:ex1} is not reproduced by the Forth-2012 files
alone.  Its required input trio is \texttt{ex1types.txt},
\texttt{ex1specs.txt}, and \texttt{ex1prog.txt}.  The accompanying artifact
contains the following paths:
\begin{center}
\small
\begin{tabular}{@{}p{0.38\linewidth}p{0.53\linewidth}@{}}
\hline
Path & Role\\
\hline
\texttt{gforth-evaluator.fs} & native checker\\
\texttt{run-evaluator-gforth.sh}, \texttt{.bat} & Unix and Windows launchers\\
\texttt{ex1types.txt}, \texttt{ex1specs.txt}, \texttt{ex1prog.txt} & complete Section~\ref{sec:ex1} profile\\
\texttt{forth2012types.txt}, \texttt{forth2012specs.txt}, \texttt{forth2012prog.txt} & larger Forth-2012-like profile\\
\texttt{testdata/positive}, \texttt{testdata/negative} & regression corpus\\
\texttt{README.md}, \texttt{EVALUATOR.md} & formats, commands, and limitations\\
\hline
\end{tabular}
\end{center}
From the artifact root, the exact command for the paper example is
\begin{lstlisting}
./run-evaluator-gforth.sh real
\end{lstlisting}
The launcher selects the three \texttt{ex1} files and creates
\texttt{ex1prog.txt.log}; that log is output, not a required input.  An archive
that omits the \texttt{ex1} trio cannot reproduce the named example.  The
essential declarations and complete source are also printed in this paper so
that the example does not depend on undocumented Forth idioms.

\subsection{Regression results}

The checked-in corpus has 29 programs: 13 positive and 16 negative.  The
positive set consists of seven Forth-2012-profile examples, five custom-profile
examples, and one article profile.  A fresh batch execution of the current
\texttt{gforth-evaluator.fs} accepts all 13 positive programs.  Fourteen of the
16 negative programs terminate with diagnostics.  The two misses are
\texttt{linear\_clash.fs} and \texttt{top\_level\_clash.fs}, both involving the
linear fragment \texttt{DP 1+}.  These cases identify gaps in top-level
linear checking.

\begin{center}
\small
\begin{tabular}{@{}p{0.34\linewidth}rrp{0.27\linewidth}@{}}
\hline
Corpus class & Programs & Expected & Current Gforth result\\
\hline
Forth-2012 positive & 7 & accept & 7 accepted\\
Custom-profile positive & 5 & accept & 5 accepted\\
Article positive & 1 & accept & 1 accepted\\
Negative & 16 & reject & 14 rejected, 2 missed\\
\hline
\end{tabular}
\end{center}

The declarative Forth-2012 environment contains 50 lines of type declarations
and 278 lines of vocabulary/control specifications.  The current native
checker is 4,424 lines of Gforth.  These counts establish nontrivial scale, but
they are not a performance benchmark or a soundness evaluation.  In
particular, the corpus is regression evidence for the implemented semantics;
it cannot expose an error shared by the checker and its specifications.

\section{Relation to Other Typed Stack Systems}
\label{sec:related}

\paragraph{Forth notation and algebraic effects.}
The Forth standard uses stack comments and numeric suffixes to distinguish
same-typed items, but does not make every comment a static judgment
\cite{forth2012,rather1996}.  Earlier work by the present author developed
algebraic composition, multiple effects, typed wildcards, inheritance, and
must-analysis operators
\cite{poial1991,poial1992,poial1994,poial2002,poial2006}.  The present artifact
puts a subset of those ideas into one source-processing tool with inferred
colon definitions, externally declared controls, and diagnostics.  Its lack of
direct multiple-effect declarations should be noted: the earlier theory is
more expressive on that point than the current input format.

\paragraph{StrongForth.}
StrongForth, developed by Stephan Becher, explores a direct response to legacy
compatibility: design and use a Forth system around static typing rather than
require an optional checker to accept arbitrary existing Forth.  Recent articles describe object-oriented typing,
programming practice, and word overloading
\cite{becher2025objects,becher2025practice,becher2025overloading}.  Overloading
is particularly relevant to words such as \texttt{0=} and \texttt{EXECUTE},
for which one weak common type can lose useful distinctions.  StrongForth is an
integrated typed dialect and can attach language meaning to those mechanisms;
the present work instead leaves the target Forth unchanged and loads an
external, currently single-effect profile.  The systems therefore answer
different questions, and StrongForth demonstrates that the restrictions of
this artifact are not restrictions of Forth-like syntax itself.

\paragraph{Optional, pluggable Forth typing.}
Riegler's 2015 thesis is the closest comparison in goals
\cite{riegler2015}.  Its Haskell prototype offers optional checkers at several
levels, from stack underflow/overflow checking to stronger static type
consistency, and treats subtyping, control structures, reference types,
multiple stack effects, compile-time programming, object orientation, and
higher-order programming.  It reports that gradual checking can improve stack
discipline with relatively few annotations, although higher-order uses need
annotations.  The current artifact contributes a native Gforth implementation,
external parser/control profiles, symbolic source traces, and the parsed-body
characterization of Section~\ref{sec:ex1}; it is narrower in multiple effects,
compile-time stack typing, and higher-order features.  These are complementary
experiments rather than evidence that one fixed strictness works for all
legacy code.

\paragraph{Compositional low-level typing.}
Stoddart and Knaggs study inference for stack-based languages
\cite{stoddart1993}, and Saabas and Uustalu give compositional systems for
stack-based low-level code \cite{saabas2006}.  Their instruction languages are
more fixed than Forth source: they do not need an open-ended outer interpreter
in which a word consumes the next name or changes parsing behavior.  In return,
a fixed instruction and control model makes conventional typing rules and
metatheory easier to state.

\paragraph{WebAssembly 3.0.}
The current online WebAssembly Core Specification identifies itself as Release
3.0, dated 23 July 2026 \cite{wasm3}.  Wasm does not infer a loop interface by
comparing one and two body executions.  A block type declares the operand
types consumed and produced; a control stack records each active block's
label type and stack height; and branches are checked against the target
label.  A loop branch targets the declared loop-label input, so every back edge
is checked directly.  Unreachable code receives a controlled form of stack
polymorphism.  The fixed bytecode and declared signatures support progress and
preservation results and mechanized validation
\cite{watt2018,mcdermott2022}.

The contrast locates the present trade-off.  Wasm gives stronger loop and
branch guarantees because its control syntax and block interfaces are fixed
and explicit.  This Forth checker tries to infer interfaces from annotation-light
source while loading parser and control behavior as data.  Symbolic identity
and retargetability are additions, but compaction, syntactic branch meet, and
the one-versus-two loop rule are costs of this particular design, not inherent
features of concatenative type checking.

\paragraph{First-class concatenative functions.}
A \emph{stack row} is the arbitrary untouched prefix introduced in Section~3.
A row variable $\rho$ lets one write the type of \texttt{DUP} schematically as
``for every prefix $\rho$ and cell type $a$,
$\rho\,a\to\rho\,a\,a$.''  A first-class quotation is itself a value carrying
such a stack transformation.  If a word accepts a quotation that must work for
\emph{every} $\rho$, the universal quantifier occurs inside the argument type.
This placement is what ``rank-$n$'' (in the common case, rank-2) polymorphism
means; it is more demanding than choosing $\rho$ once for the whole call.

A recent account of the IV language avoids automatic inference of these nested
quantifiers by requiring definition annotations, extending operator types
across stack depths, and giving execution and composition operators explicit
parameters \cite{dinikeev2025}.  The current checker has indexed cell variables
and an implicit frame, but no type for quotation values and no algorithm that
unifies explicit row variables.  Supporting first-class quotations therefore
requires an extension of the type language, not merely another vocabulary
entry.

\section{Limitations and Research Directions}

The implementation realizes a restricted static analysis with the following
scope limitations.
\begin{itemize}
  \item Acceptance is relative to trusted primitive and parser specifications.
        An incorrect declaration can make the result unsound with respect to
        concrete execution.
  \item Comparable type meet rejects incomparable types even if the declared
        poset contains a useful third lower bound.  The input format also has
        no direct multiple-effect overload for one word.
  \item Compaction stores away some implicit correlations instead of limiting
        abbreviation to Forth-style output.
  \item Branch meet is syntactic and may strengthen a precondition rather than
        compute a safe summary of both paths.  \texttt{MAYSWAP} is a concrete
        example.  In addition, the native unequal-depth helper currently
        aligns its residual comparison at offset zero rather than $k$.
  \item $\pi^\star_2$ is not an inductive loop-invariant computation; the
        \texttt{P} example in Section~\ref{sec:loops} is accepted with a
        non-idempotent reported summary.
  \item Control words are parsed structurally, but their standard compile-time
        control-flow-stack effects are not typed.
  \item Non-local effects such as \texttt{THROW}, return-stack manipulation,
        dynamic execution, and unrestricted metaprogramming are approximated
        or outside the model.
  \item The parsed-body theorem excludes recursive dictionary construction and
        forward-reference machinery.
  \item Two checked-in top-level negative tests are not currently rejected.
\end{itemize}

The immediate engineering priorities are to retain identity classes internally,
format inferred effects in consecutive Forth-standard style, correct the
unequal-depth alignment, and repair the top-level clash path.  A union--find
representation would avoid repeated whole-effect substitution.  Language-level
extensions include multiple primitive effects, an explicit compile-time
control-flow stack, row variables and quotation types, disjunctive branch
summaries, and a control-flow fixpoint for loops.  A semantic development would
add an operational model for primitives, variance-aware pre/post contracts,
and a mechanized soundness argument for the parser-to-equation translation.

\section{Conclusion}

This work is specifically a configurable source-level checker for Forth, not a
general account of all concatenative type systems.  Its useful core is small:
subtype-aware boundary unification composes straight-line words, indexed cells
carry value correlations through stack shuffles, and external profiles describe
Forth's parser, dictionary words, and control syntax.  Once a non-recursive body
has been parsed, the checker performs a deterministic bottom-up evaluation of
those declarations.

Branches and loops are where stack depth alone stops being enough.  The paper
has made the current answers explicit rather than treating the names
\texttt{glb} and \texttt{pistar} as guarantees: branch combination is a partial
syntactic contract with known alignment and generality problems, and repetition
compares only one execution with two.  The \texttt{MAYSWAP} and \texttt{P}
examples delimit those claims.  StrongForth, Riegler's pluggable checker, and
WebAssembly demonstrate other viable points in the design space.

The artifact is therefore useful as configurable, executable stack-effect
documentation and as an experimental platform, but not yet as a sound static
semantics for arbitrary Forth.  Preserving identities internally, adopting
standard-style output, supporting multiple effects and compile-time control
items, and replacing local repetition by an inductive analysis are concrete
steps from this prototype toward that stronger goal.

\begin{thebibliography}{99}

\bibitem{becher2025objects}
S. Becher.
\newblock Statische Typpr\"ufung und Objektorientierung.
\newblock \emph{Forth--Magazin Vierte Dimension}, 41(1):11--14, 2025.
\newblock \url{https://forth-ev.de/wiki/res/lib/exe/fetch.php/vd-archiv:4d2025-01.pdf}.

\bibitem{becher2025practice}
S. Becher.
\newblock Statische Typpr\"ufung: Programmierpraxis.
\newblock \emph{Forth--Magazin Vierte Dimension}, 41(2):15--20, 2025.
\newblock \url{https://forth-ev.de/wiki/res/lib/exe/fetch.php/vd-archiv:4d2025-02.pdf}.

\bibitem{becher2025overloading}
S. Becher.
\newblock Statische Typpr\"ufung: \"Uberladung von W\"ortern.
\newblock \emph{Forth--Magazin Vierte Dimension}, 41(3):14--16, 2025.
\newblock \url{https://forth-ev.de/wiki/res/lib/exe/fetch.php/vd-archiv:4d2025-03.pdf}.

\bibitem{cousot1977}
P. Cousot and R. Cousot.
\newblock Abstract interpretation: A unified lattice model for static analysis
of programs by construction or approximation of fixpoints.
\newblock In \emph{POPL}, pp. 238--252, 1977.
\newblock \href{https://doi.org/10.1145/512950.512973}{doi:10.1145/512950.512973}.

\bibitem{dinikeev2025}
A. M. Dinikeev.
\newblock Type system for a statically typed concatenative programming
language with first class function support.
\newblock \emph{Information-measuring and Control Systems}, no. 5,
pp. 15--25, 2025.
\newblock \href{https://doi.org/10.18127/j20700814-202505-02}{doi:10.18127/j20700814-202505-02}.

\bibitem{forth2012}
Forth 200x Standardisation Committee.
\newblock \emph{Forth 2012 Standard}.
\newblock Draft proposed standard, 10 November 2014.
\newblock \url{https://forth-standard.org/standard/intro}.

\bibitem{mcdermott2022}
D. McDermott, Y. Morita, and T. Uustalu.
\newblock A type system with subtyping for WebAssembly's stack polymorphism.
\newblock In \emph{ICTAC 2022}, LNCS 13572, pp. 305--323, 2022.
\newblock \href{https://doi.org/10.1007/978-3-031-17715-6_20}{doi:10.1007/978-3-031-17715-6\_20}.

\bibitem{poial1991}
J. Pöial.
\newblock Algebraic specifications of stack effects for Forth programs.
\newblock In \emph{1990 FORML Conference Proceedings, EuroFORML'90},
pp. 282--290, 1991.

\bibitem{poial1992}
J. Pöial.
\newblock Multiple stack-effects of Forth programs.
\newblock In \emph{1991 FORML Conference Proceedings, euroFORML'91},
pp. 400--406, 1992.

\bibitem{poial1994}
J. Pöial.
\newblock Algebraic specifications of stack effects.
\newblock \emph{Forth Dimensions}, 16(4):18--20,
November/December 1994.
\newblock \url{https://www.forth.org/fd/FD-V16N4.pdf}.

\bibitem{poial2002}
J. Pöial.
\newblock Stack effect calculus with typed wildcards, polymorphism and
inheritance.
\newblock Abstract in \emph{18th EuroForth Conference}, p. 38, 2002;
accompanying presentation slides.
\newblock \url{https://www.complang.tuwien.ac.at/anton/euroforth/ef02/poial02.pdf}.

\bibitem{poial2006}
J. Pöial.
\newblock Typing tools for typeless stack languages.
\newblock In \emph{Proceedings of the 22nd EuroForth Conference},
Cambridge, England, pp. 40--46, 2006.
\newblock \url{https://www.complang.tuwien.ac.at/anton/euroforth/ef06/poial06.pdf}.

\bibitem{poial2008}
J. Pöial.
\newblock Java framework for static analysis of Forth programs.
\newblock In \emph{Proceedings of the 24th EuroForth Conference},
Vienna, Austria, pp. 20--23, 2008.
\newblock \url{https://www.complang.tuwien.ac.at/anton/euroforth/ef08/papers/poial.pdf}.

\bibitem{riegler2015}
G. Riegler.
\newblock \emph{Evaluation and Implementation of an Optional, Pluggable Type
System for Forth}.
\newblock Diploma thesis, Technische Universit\"at Wien, 2015.
\newblock \href{https://doi.org/10.34726/hss.2015.27899}{doi:10.34726/hss.2015.27899}.
\newblock \url{https://repositum.tuwien.at/handle/20.500.12708/7117}.

\bibitem{rather1996}
E. D. Rather, D. R. Colburn, and C. H. Moore.
\newblock The evolution of Forth.
\newblock In \emph{History of Programming Languages---II}, pp. 625--670,
ACM, 1996.
\newblock \href{https://doi.org/10.1145/234286.1057832}{doi:10.1145/234286.1057832}.

\bibitem{saabas2006}
A. Saabas and T. Uustalu.
\newblock Compositional type systems for stack-based low-level languages.
\newblock In J. Gudmundsson and C. B. Jay, editors,
\emph{Proceedings of CATS 2006}, CRPIT 51, pp. 27--39.
Australian Computer Society, 2006.
\newblock \url{https://cs.ioc.ee/~tarmo/papers/saabas-uustalu-cats06.pdf}.

\bibitem{stoddart1993}
B. Stoddart and P. J. Knaggs.
\newblock Type inference in stack based languages.
\newblock \emph{Formal Aspects of Computing}, 5(4):289--298, 1993.
\newblock \href{https://doi.org/10.1007/BF01212404}{doi:10.1007/BF01212404}.

\bibitem{wasm3}
WebAssembly Community Group.
\newblock \emph{WebAssembly Core Specification, Release 3.0}.
\newblock Editor: A. Rossberg, 23 July 2026.
\newblock \url{https://webassembly.github.io/spec/core/}.

\bibitem{watt2018}
C. Watt.
\newblock Mechanising and verifying the WebAssembly specification.
\newblock In \emph{CPP 2018}, pp. 53--65, ACM, 2018.
\newblock \href{https://doi.org/10.1145/3167082}{doi:10.1145/3167082}.

\end{thebibliography}

\end{document}
